% !PRINTED 466
\weberpage{466}

%% [running head, not body text -- commented out 2026-08-07] Page furniture (running
%% head, folio, signature) is excluded from this edition by standing policy.
%% \noindent 466 \hfill Nineteenth Chapter. \hfill \S\,126.

This procedure is not applicable to the case $m=1$, because $\Phi_n(u,u)$ vanishes identically for $n=1$, $\omega=i$ (but also only for this value); for $m=1$, however, we have already found earlier (\S\,83):
\[
f(i)^{24}=64,
\]
so that in this case too the theorem stated holds.

If $m$ is odd, then for odd $A$, $C$ the middle coefficient $B$ is divisible by 4; if, on the other hand, $m$ is even, then under the same assumption $B$ is not divisible by 4. But by the substitution $\omega+1$ for $\omega$ we achieve, in the case of an even $m$ as well, that $B$ becomes divisible by 4; but then $C$ also becomes even. By the interchange $(\omega,\omega+1)$, however, $f(\omega)^{24}$ passes into $-f_1(\omega)^{24}$, so that we can also give our theorem the following expression:

2. \gesperrt{If $\omega$ is the root of the quadratic equation:}
\[
A\omega^2+B\omega+C=0,
\]
\gesperrt{in which $A$ is odd and $B$ is even, then, according as $D\equiv 4$ or $D\equiv 0\pmod{8}$, $f(\omega)^{24}$ or $f_1(\omega)^{24}$ is a class invariant.}

Since $A$ and $C$ cannot both be even, only three cases are possible, of which, by the interchanges
\[
(\omega,\omega+1),\quad \left(\omega,-\frac{1}{\omega}\right)
\]
the second is reduced to the first and the third to the second. In these cases, then, we have:
\weq{13}{
\begin{array}{lll}
A\equiv 1, & C\equiv 1\ (\mathrm{mod}\ 2), & \text{class invariant } f(\omega)^{24},\\
A\equiv 1, & C\equiv 0 & \text{\glqq}\quad\text{\glqq}\quad f_1(\omega)^{24},\\
A\equiv 0, & C\equiv 1 & \text{\glqq}\quad\text{\glqq}\quad f_2(\omega)^{24}.
\end{array}
}

In particular, therefore, if we take as our basis the principal form of the discriminant $-4m$, $(1,0,m)$, then $f(\sqrt{-m})^{24}$ for odd, and $f_1(\sqrt{-m})^{24}$ for even $m$, is a class invariant. These two numbers have a real positive value, and in the calculations that follow further below they shall be taken into consideration by preference.

From equation (1) we conclude (Bd.\,II, \S\,154, 11.), since $j(\omega)$ is an integral algebraic number, that $f(\omega)^{24}$ must also be an integral algebraic number, and the form of equation (1) shows that the norm of this number is a power of 2.
